\documentclass[10pt,a4paper]{amsart}
\usepackage[margin=19mm]{geometry}
\usepackage{amsmath,amssymb,mathtools}
\usepackage[scr=rsfs,cal=euler]{mathalfa}
\usepackage{xfrac}
\usepackage[unicode,hidelinks,bookmarks=false]{hyperref}
\newcommand{\ie}{i.e.\ }
\newcommand{\cf}{cf.\ }
\newcommand{\resp}{respectively\ }
\newcommand{\Subsection}{\subsection}
\providecommand{\nobreakdash}{\nobreak\hbox{-}}
\setlength{\emergencystretch}{2em}
\multlinegap0pt
\usepackage{amssymb}
\usepackage{tikz}
\newtheorem{lemma}{Lemma}
\DeclareMathOperator{\rht}{ht}
\title{Recurrences for permutations with long increasing subsequences}
\author{Manuel Kauers, Chen Wang}
\date{Source: arXiv:2609.02220v1 (CC BY 4.0)}
\begin{document}
\maketitle

\noindent\emph{Source and licence.} Recurrences for permutations with long increasing subsequences. Version arXiv:2609.02220v1 (CC BY 4.0), distributed by arXiv under the Creative Commons Attribution 4.0 International licence, http://creativecommons.org/licenses/by/4.0/, as recorded in the arXiv licence metadata for this exact version and shown on the abstract page https://arxiv.org/abs/2609.02220v1. Full licence notice: \texttt{licenses/arxiv-2609.02220-CC-BY-4.0.txt}.\par\smallskip

\noindent\textit{Editorial scope.} Counted unit: the article's lemma lem:parseval (source0.tex lines 278--283). The article's complete original proof of it is delivered with it (source0.tex lines 285--329, one \texttt{proof} environment of 45 lines). No mathematics is written, repaired or re-derived: the statement and the proof are the article's own lines, in the article's own notation, and no retained line is substituted or re-flowed.  No further source-local context is needed: every cross-reference of every retained line resolves inside the two delivered spans.  Nothing was omitted from any retained span, so the package carries no omission record.  No retained line contains a citation, so the wrapper carries no bibliography and no bibliography entry.  No retained line contains a figure environment, an image-inclusion command or drawing code, so no image is delivered and the package declares no asset dependency.  Editorial formatting only, in addition: an AMS \texttt{amsart} preamble that restates the article's own declarations and macros used by the retained lines, namely the environments \texttt{lemma} on separate counters, together with the standard packages those lines need.  The retained environments print in this document as Lemma 1 in order of appearance, whereas the article prints the same items with the numbering of its own full document.  The complete native source of the article as distributed in the arXiv e-print for this exact version is bundled unchanged as \texttt{sources/209183/source0.tex}.
\begin{lemma}\label{lem:parseval}
For every $V\in P_m$,
\begin{equation}\label{eq:parseval}
    \sum_{s=1}^m\|D_sV\|^2=m\|V\|^2.
\end{equation}
\end{lemma}

\begin{proof}
We write
\[
V=\sum_{\lambda\vdash m}v_\lambda x_{\lambda}.
\]
Standard character orthogonality implies
\begin{equation}\label{eq:vnorm}
\|V\|^2=\frac{1}{m!}\sum_{\sigma\in S_m}\left(\sum_{\lambda\vdash m}v_\lambda\chi_{\lambda}(\sigma)\right)^2.
\end{equation}

On the other hand, the Murnaghan--Nakayama rule gives
\begin{equation*}
    D_sV=\sum_{\lambda\vdash m}v_\lambda\sum_{\substack{\mu\vdash m-s\\
\lambda\setminus\mu\text{ is a rim hook}}}
(-1)^{\rht(\lambda\setminus\mu)}x_\mu,
\end{equation*}
and therefore
\begin{align}
    \|D_sV\|^2
    &=\frac{1}{(m-s)!}\sum_{\tau\in S_{m-s}}
      \left(\sum_{\substack{\lambda\vdash m\\ \mu\vdash m-s\\
      \lambda\setminus\mu\text{ is a rim hook}}}
      (-1)^{\rht(\lambda\setminus\mu)}v_\lambda\chi_\mu(\tau)\right)^2\nonumber\\
    &=\frac{1}{(m-s)!}\sum_{\tau\in S_{m-s}}
      \left(\sum_{\lambda\vdash m}v_\lambda\chi_\lambda(c_s\tau)\right)^2,
      \label{eq:dvnorm}
\end{align}
where $c_s$ is an $s$-cycle disjoint from $\tau$.

Combining \eqref{eq:vnorm} and \eqref{eq:dvnorm}, we see that
\eqref{eq:parseval} is equivalent to
\begin{equation*}
    \sum_{\sigma\in S_m}\left(\sum_{\lambda\vdash m}v_\lambda\chi_{\lambda}(\sigma)\right)^2
    =\sum_{s=1}^{m}\frac{(m-1)!}{(m-s)!}\sum_{\tau\in S_{m-s}}
      \left(\sum_{\lambda\vdash m}v_\lambda\chi_\lambda(c_s\tau)\right)^2.
\end{equation*}
This last identity can be proved bijectively.  Let $S_m$ act on
$\{1,2,\dots,m\}$.  There are exactly
\[
 \binom{m-1}{s-1}(s-1)! = \frac{(m-1)!}{(m-s)!}
\]
possible $s$-cycles containing $m$, and every permutation is uniquely
the product of its cycle containing $m$ and a permutation on the
remaining points.
\end{proof}

\end{document}
